\clearpage\chapter{Bass--Serre theory}
\section*{1. Overview}
\subsection*{Intuitive}
\paragraph{The problem}

A group can be difficult to understand when its elements are presented by many relations. Bass--Serre theory gives a geometric picture: make the group act on a tree, a connected graph with no loops. Vertices and edges then carry subgroups, and the way those subgroups are attached records whether the group is a free product, an amalgamated product, or an HNN extension.



\subsection*{Formal}
Let a group $G$ act without inversions on a tree $T$. Each vertex $v$ and edge $e$ has a stabilizer subgroup $G_v$ or $G_e$. The quotient graph $G\backslash T$, together with these stabilizers and the inclusion maps from edge groups to vertex groups, is a graph of groups. Bass--Serre theory reconstructs $G$ as the fundamental group of this graph of groups and uses the action on $T$ to study its structure.
\section*{2. History}
\subsection*{Historical development}
\begin{historylist}
\paragraph{History and motivation}

Jean-Pierre Serre developed the theory in the 1970s, with Hyman Bass making major contributions to its formalization. Serre's book \emph{Trees} grew partly from the study of algebraic groups whose Bruhat--Tits buildings are trees. The method quickly became central to geometric group theory and three-manifold topology.

The older algebraic constructions were free products with amalgamation and HNN extensions. Bass--Serre theory did not replace them; it explained them geometrically through covering spaces, fundamental groups, and group actions. A graph of groups is a one-dimensional analogue of an orbifold: groups sit at vertices and edges, and inclusions describe how local symmetry is glued into global symmetry \cite{bass_serre_theory_ref1,bass_serre_theory_ref3}.

\end{historylist}
\section*{3. Theory}
\subsection*{Core theory}
\paragraph{Trees and group actions}

A tree is a graph in which any two vertices are joined by exactly one simple path. Let a group $G$ act on a tree $X$ by graph automorphisms. The stabilizer of a vertex $v$ is
\[
 G_v=\{g\in G:gv=v\},
\]
and the stabilizer of an edge $e$ is defined similarly. We assume first that there are no edge inversions: no element flips an edge while fixing its midpoint.

The quotient graph $X/G$ records the orbits of vertices and edges. The quotient alone loses information, so attach the vertex stabilizer to each quotient vertex and the edge stabilizer to each quotient edge. The inclusion of an edge stabilizer into the stabilizer of its endpoint gives the attaching maps. The result is a graph of groups.

The action is geometric information about $G$. If $G$ acts freely on a tree, no nonidentity element fixes a vertex, and the quotient graph has fundamental group isomorphic to $G$. This immediately implies that a group acting freely on a tree is free.

\paragraph{Graphs of groups and their fundamental group}

Choose a maximal subtree $T$ in the underlying graph of a graph of groups $\mathcal A$. The fundamental group is generated by all vertex groups and one stable letter for each edge outside $T$, subject to the internal relations in every vertex group, the relations identifying the two images of each edge-group element along an edge in $T$, and conjugacy relations involving the stable letters for edges outside $T$.

If the graph has two vertices joined by one edge, with vertex groups $A$ and $B$ and edge group $C$ embedded in each, the fundamental group is the amalgamated product
\[
 A*_C B.
\]
If the graph has one vertex and one loop edge with edge group $C$ embedded twice in $A$, the fundamental group is an HNN extension
\[
 \langle A,t\mid t^{-1}\alpha(c)t=\beta(c)\text{ for }c\in C\rangle.
\]
Free products are the special case $C$ trivial.

\paragraph{Bass--Serre covering trees}

Every graph of groups $\mathcal A$ with a chosen base vertex has a canonical covering tree $\widetilde{\mathcal A}$. Its vertices can be represented by cosets of vertex groups in the fundamental group, and its edges by cosets of edge groups. The fundamental group acts on this tree without edge inversions. The quotient graph of groups is the original $\mathcal A$.

Conversely, if $G$ acts on a tree $X$ without inversions, the quotient graph of groups has a covering tree naturally isomorphic to $X$. The construction generalizes the universal cover of an ordinary graph. Ordinary covering theory has a fundamental group acting on a tree; Bass--Serre theory allows nontrivial stabilizers, so the local groups become part of the cover.

\paragraph{The fundamental theorem}

Let $G$ act on a tree $X$ without inversions, let $\mathcal A$ be the quotient graph of groups, and choose a base vertex $v$. Then
\[
 G\cong\pi_1(\mathcal A,v),
\]
and $X$ is equivariantly isomorphic to the Bass--Serre covering tree of $\mathcal A$. The isomorphisms respect both the group action and the graph structure.

This theorem is powerful in both directions. To understand an action, pass to the quotient graph of groups and obtain a group decomposition. To build an action for a group given by an amalgam or HNN extension, build its Bass--Serre tree.

\paragraph{Basic consequences}

If a graph of groups has a spanning tree and fundamental group $G$, the natural maps from its vertex groups into $G$ are injective. A finite-order element of $G$ is conjugate to an element of a vertex group. More generally, every finite subgroup is conjugate into some vertex group. If the quotient graph is finite and all vertex groups are finite, $G$ is virtually free: it contains a free subgroup of finite index.

Finite graphs of finitely generated vertex groups give finitely generated fundamental groups. If vertex groups are finitely presented and edge groups finitely generated, the resulting group is finitely presented. These statements make a decomposition useful for transferring finiteness properties from local groups to the global group.

\paragraph{Trivial and nontrivial actions}

An action is trivial if the whole group fixes a vertex. A graph of groups is correspondingly trivial when its underlying graph is a tree and the edge inclusions collapse the other vertex groups into the chosen base vertex group. Such a decomposition does not reveal a genuine splitting.

Nontrivial actions have no global fixed vertex and encode a real group splitting. The distinction matters in applications: rooted-tree actions may be interesting even when they fix a root, but Bass--Serre decomposition arguments generally seek actions that cannot be reduced to one vertex group.

\paragraph{Ends, property T, and accessibility}

Stallings' theorem says that a finitely generated group has more than one end exactly when it admits a nontrivial splitting over a finite subgroup. In Bass--Serre language, the group acts nontrivially on a tree with finite edge stabilizers. This turns a coarse geometric property of the Cayley graph into an algebraic decomposition.

If a group has Kazhdan's property (T), every action on a tree without inversions has a global fixed vertex. Such a group has Serre's property FA and cannot split nontrivially in the Bass--Serre sense. This provides an obstruction to decomposition.

Accessibility theorems bound the complexity of splittings. Dunwoody proved that a finitely presented group has a bound on the number of edges in reduced splittings over finite subgroups. Bestvina and Feighn generalized accessibility to splittings over small subgroups. These results say that decomposition cannot continue indefinitely once finiteness hypotheses are imposed \cite{bass_serre_theory_ref4,bass_serre_theory_ref5}.

\section*{4. Important theorems and results}
\begin{enumerate}[leftmargin=*,itemsep=0.35em]

\item \textbf{Fundamental theorem of Bass--Serre theory.} A group acting without inversions on a tree is the fundamental group of its quotient graph of groups, and the tree is the associated Bass--Serre covering tree.

\item \textbf{Free-action theorem.} A group acting freely on a tree is a free group.

\item \textbf{Kurosh subgroup theorem.} Subgroups of a free product decompose as a free product of conjugates of subgroups of the factors together with a free group.

\item \textbf{Stallings' ends theorem.} A finitely generated group has more than one end exactly when it splits nontrivially over a finite subgroup.

\item \textbf{Property FA consequence.} A group with Kazhdan's property (T) fixes a vertex in every tree action without inversions and therefore has no nontrivial Bass--Serre splitting.

\item \textbf{Accessibility.} Under suitable finite-presentation assumptions, the complexity of reduced splittings is bounded.
\end{enumerate}
\noindent\emph{Source for these results:} \cite{bass_serre_theory_ref1}.

\section*{5. Applications}
Bass--Serre theory applies group actions on trees to groups assembled from simpler pieces. The resulting tree records free products, amalgamations, and HNN extensions in a form that exposes their subgroup structure.
\begin{enumerate}[leftmargin=*,itemsep=0.35em]
\item \textbf{Free product.} The group $A*B$ acts on a tree whose vertices are cosets of $A$ and $B$ and whose edges connect compatible cosets. The quotient is one edge with two vertices and trivial edge group. The tree makes reduced words alternate between factors.
\item \textbf{Amalgamated product.} If two groups share a subgroup $C$, the group $A*_C B$ glues the two vertex groups along $C$. The edge stabilizer records exactly the elements identified in the gluing. This appears in three-manifold groups and combination theorems.
\item \textbf{HNN extension.} An HNN extension joins two copies of a subgroup inside one group with a stable letter. The Bass--Serre tree records repeated conjugation of one subgroup into another. It is a natural model for groups built by adding a symmetry that identifies two substructures.
\item \textbf{Three-manifolds and geometry.} Splittings over boundary or torus subgroups help analyze fundamental groups of three-manifolds. The tree records how pieces are assembled and interacts with geometric decompositions.
\item \textbf{Geometric group theory.} Actions on trees are one-dimensional models of nonpositive-curvature actions. They help study subgroup structure, free subgroups, ends, and algorithmic questions about group presentations.
\end{enumerate}


\theoryconnections{bass_serre_theory}{%
\item Group theory supplies the symmetries, graph theory supplies vertices and edges, and topology explains why a graph of groups behaves like a space assembled from pieces.
\item The fundamental group of a graph records the possible ways to travel around it; when groups are attached to vertices and edge groups describe the overlaps, the resulting action on a tree records how the whole group splits.
\item An amalgamated product identifies a shared subgroup, while an HNN extension identifies two copies of one subgroup \cite{bass_serre_theory_ref1,bass_serre_theory_ref2}.
}{%
\item The action on a tree turns an algebraic splitting into geometry: vertex stabilizers are the pieces, edge stabilizers are the overlaps, and the quotient graph displays the assembly plan.
\item This makes free products and subgroup structure visible, gives criteria for decomposing fundamental groups of three-manifolds, and supplies geometric tools for studying groups that are too large to classify by generators and relations alone \cite{bass_serre_theory_ref1,bass_serre_theory_ref4}.
}

\section*{Glossary}
\begin{description}[leftmargin=*,style=nextline,itemsep=0.3em]

\item[\textbf{Tree.}] A connected graph with no cycles; in Bass--Serre theory, a group acts on a tree and the structure of that action encodes how the group decomposes.

\item[\textbf{Graph of groups.}] A graph in which each vertex and edge is labeled by a group, with inclusion maps from edge groups to vertex groups; its fundamental group assembles the local groups into a global group.

\item[\textbf{Amalgamated product.}] A group formed by taking two groups $A$ and $B$ that share a common subgroup $C$ and identifying $C$ inside both, written $A*_C B$.

\item[\textbf{HNN extension.}] A group constructed from a single group $A$ and a new element $t$ (the stable letter) that conjugates one subgroup of $A$ into another.

\item[\textbf{Stabilizer.}] The subgroup of elements in a group action that fix a given vertex or edge; stabilizers record the local symmetry at each point of the tree.

\item[\textbf{Ends of a group.}] The number of ways the Cayley graph of a group extends to infinity after removing any finite region; more than one end signals a nontrivial splitting.

\item[\textbf{Free product.}] The amalgamated product $A*_C B$ when the common subgroup $C$ is trivial, producing the most general way to combine two groups without identifying any elements.

\item[\textbf{Covering tree.}] The canonical Bass--Serre tree associated with a graph of groups, whose vertices and edges are cosets of vertex and edge groups.

\item[\textbf{Quotient graph.}] The graph obtained by collapsing each orbit of vertices and edges under a group action.

\item[\textbf{Kazhdan's property (T).}] A group-theoretic property meaning every action on a tree without inversions has a global fixed vertex.

\item[\textbf{Virtually free.}] A group containing a free subgroup of finite index.

\end{description}

\section*{7. Sample Questions}
\emph{Exercise reference:} \cite{bass_serre_theory_ref1}. The cited edition is the source for the exercise set; consult its Exercises section for the official statement and numbering.
\begin{enumerate}[leftmargin=*,itemsep=0.9em]

\item \textbf{Exercise 1.} Let $G=\mathbb{Z}/2\mathbb{Z}*\mathbb{Z}/3\mathbb{Z}$ be the free product of cyclic groups of orders 2 and 3. Write down the graph of groups for this decomposition.

\emph{Solution.} The graph has two vertices, one labeled $\mathbb{Z}/2\mathbb{Z}$ and one labeled $\mathbb{Z}/3\mathbb{Z}$, joined by a single edge whose edge group is trivial. The fundamental group of this graph of groups is $\mathbb{Z}/2\mathbb{Z}*\mathbb{Z}/3\mathbb{Z}$. The Bass--Serre tree is a tree whose vertices alternate between cosets of $\mathbb{Z}/2\mathbb{Z}$ and cosets of $\mathbb{Z}/3\mathbb{Z}$.

\item \textbf{Exercise 2.} The group $\mathbb{Z}$ acts freely on the real line $\mathbb{R}$ (a tree) by $n\cdot x=x+n$. Identify the quotient graph of groups and verify the fundamental theorem.

\emph{Solution.} The quotient $\mathbb{R}/\mathbb{Z}$ is a single vertex with a single loop edge. Since the action is free, all stabilizers are trivial. The graph of groups has one vertex with trivial group and one loop edge with trivial group. By the fundamental theorem, $G\cong\pi_1(\mathcal{A},v)$. The fundamental group of a graph with one vertex and one edge is $\mathbb{Z}$, confirming $G\cong\mathbb{Z}$.

\item \textbf{Exercise 3.} Let $A=\mathbb{Z}/4\mathbb{Z}$, $B=\mathbb{Z}/6\mathbb{Z}$, and $C=\mathbb{Z}/2\mathbb{Z}$ embedded via the unique injections $C\hookrightarrow A$ and $C\hookrightarrow B$. Write the presentation of the amalgamated product $A*_C B$.

\emph{Solution.} Let $A=\langle a\mid a^4=1\rangle$, $B=\langle b\mid b^6=1\rangle$, and $C=\langle c\mid c^2=1\rangle$ with $C\hookrightarrow A$ sending $c\mapsto a^2$ and $C\hookrightarrow B$ sending $c\mapsto b^3$. The amalgamated product is $A*_C B=\langle a,b\mid a^4=1,\,b^6=1,\,a^2=b^3\rangle$.

\item \textbf{Exercise 4.} Show that the fundamental group of the figure-eight graph (one vertex, two loop edges $e_1,e_2$) with trivial vertex and edge groups is the free group $F_2$.

\emph{Solution.} Choose the single vertex as the spanning tree (which is trivially a tree). Both edges $e_1,e_2$ are outside the spanning tree, each contributing a stable letter. The fundamental group is generated by the trivial vertex group together with stable letters $t_1,t_2$ for $e_1,e_2$, with no relations beyond those in the vertex group. Hence $\pi_1(\mathcal{A})=\langle t_1,t_2\rangle\cong F_2$, the free group on two generators.

\item \textbf{Exercise 5.} Let $F$ be a free group of rank 2 and $H$ the subgroup generated by $a^2$ and $b^2$ where $a,b$ are the free generators. Use the Kurosh subgroup theorem to describe $H$.

\emph{Solution.} The Kurosh subgroup theorem states that a subgroup of a free product decomposes as a free product of conjugates of subgroups of the factors, together with a free group. Since $F_2=\mathbb{Z}*\mathbb{Z}$ and $H\le F_2$ is itself free (every subgroup of a free group is free), the theorem says $H$ is a free group. The rank of $H$ can be computed via the Schreier index formula: $H$ has index $[F_2:H]=\infty$ (since $\langle a^2,b^2\rangle$ has infinite index), so $H$ is free of infinite rank.

\end{enumerate}
\section*{8. References}
\begin{thebibliography}{99}
\bibitem{bass_serre_theory_ref1} J.-P. Serre, \emph{Trees}, Springer, 1980.
\bibitem{bass_serre_theory_ref2} H. Bass, "Covering theory for graphs of groups," \emph{Journal of Pure and Applied Algebra}, 1993.
\bibitem{bass_serre_theory_ref3} A. Hatcher, \emph{Algebraic Topology}, Cambridge University Press, 2002.
\bibitem{bass_serre_theory_ref4} W. Dicks and M. J. Dunwoody, \emph{Groups Acting on Graphs}, Cambridge University Press, 1989.
\bibitem{bass_serre_theory_ref5} D. E. Cohen, \emph{Combinatorial Group Theory: A Topological Approach}, Cambridge University Press, 1989.
\end{thebibliography}
